\subsection{OPS và context đầu bộ nhân}

OPS đếm bit 1 trên phần fraction 12 bit của A và B. Mỗi phép
đếm được chia thành hai nhóm 6 bit, rồi cộng hai số đếm cục bộ.
OPS1 chọn B làm X chỉ khi count B nhỏ hơn count A. Trường hợp
hòa giữ A làm X. Do X quyết định số token cần phát, OPS1 có thể
giảm số vòng cho cùng cặp A/B. OPS0 giúp kiểm vai cố định.
Y gồm bit ẩn và phần fraction 12 bit, nên bus Y có 13 bit.
Dấu tích là XOR, số mũ cơ sở là $sf_A+sf_B$ signed 11 bit.

\rtlfig{OPS, SAC và SBM trong bộ nhân FULL\_TERM26/Q24. T1/T2 chốt token; T3 chốt tổng. Hai đường hồi tiếp phục vụ cập nhật mask và cộng term.}{fig:multiplier-stages}{
\node[pin] (ab) at (-1.55,.46) {$f_A[11:0], f_B[11:0]$};
\node[pin] (cfg) at (5.4,.46) {\sig{ops\_q}, \sig{n\_q[3:0]}};
\node[unit,text width=74mm] (ops) at (-1.55,-1.06) {\hw{OPS}\\
 mỗi toán hạng: 2 đếm 6 bit $\to$ tổng 4 bit\\
 B làm X nếu OPS1 và $\mathrm{popcount}(B)<\mathrm{popcount}(A)$};
\draw[wire] (ab)--(ops);
\draw[enable] (cfg.south)|-node[tag,above]{OPS}(ops.east);
\node[ff,text width=64mm] (sacreg) at (-2.45,-3.02) {\hw{Context SAC tại launch}\\
 \sig{fx\_q[11:0]} $\leftarrow X$; scale/count $\leftarrow0$\\
 \sig{limit\_q[3:0]} $\leftarrow n$};
\node[ff,text width=47mm] (y) at (4.55,-3.02) {\hw{Context Y và metadata}\\\sig{y\_q[12:0]} $\leftarrow \{1,f_Y\}$\\dấu, $sf$, \sig{input\_cut}};
\draw[wire] (ops.south -| sacreg)--node[tag,right]{X:12}(sacreg);
\draw[wire] (ops.south)--++(0,-.31)-|(y.north) node[pos=.25,tag,above]{Y:13};
\node[unit,text width=64mm] (sac) at (-2.45,-5.18) {\hw{SAC + điều khiển phát token}\\
 LOD12 $\to (K,v)$; $S' = S+K+1$\\
 $F'=\{F[10:0],0\}\ll K$ khi $v=1$\\
 hết bit nếu $F\mathbin{\&}(F-1)=0$};
\draw[wire] (sacreg)--node[tag,right]{$F,S$, count, limit}(sac);
\draw[wire] (sac.west)--++(-.52,0)|-(sacreg.west);
\node[explain,text width=46mm] at (4.55,-5.2) {\hw{Dừng phát}\\init nếu $F=0$ hoặc $n=0$.\\
 last nếu init, hết bit,\\hoặc count+1 đạt $n$.};
\node[ff,text width=64mm] (t1) at (-2.45,-7.25) {\hw{T1} \quad scale 4, iteration 4\\valid, first, last, init, cut};
\draw[wire] (sac)--(t1);
\node[unit,text width=45mm] (base) at (4.55,-7.25) {\hw{Cơ số Q24}\\$B_0=\{0,\sig{y\_q},12'b0\}$};
\draw[wire] (y.east)--++(.33,0)|-(base.east);
\node[unit,text width=64mm] (sbm) at (-2.45,-8.98) {\hw{SBM: dịch phải 26 bit}\\
 các cấp 1/2/4/8; term $=B_0\gg S'$\\init-only chọn term $=0$};
\draw[wire] (t1)--(sbm);
\draw[wire] (base.south)--++(0,-.35)|-(sbm.east) node[pos=.63,tag,above]{26};
\node[ff,text width=64mm] (t2) at (-2.45,-10.78) {\hw{T2} \quad term 26, iteration 4\\valid, first, last, cut};
\draw[wire] (sbm)--(t2);
\node[pick,text width=32mm] (m) at (4.55,-11.04) {first ? $B_0$ : acc};
\draw[wire] (base)--(m);
\node[unit,text width=64mm] (sum) at (-2.45,-12.53) {\hw{Cộng mở rộng 27 bit}\\base + term; bit 26 kiểm carry};
\draw[wire] (t2)--node[tag,right]{term:26}(sum);
\draw[wire] (m.south)|-(sum.east) node[pos=.62,tag,above]{base:26};
\draw[enable] (t2.east)|-(m.west) node[midway,tag,above]{first};
\node[ff,text width=64mm] (t3) at (-2.45,-14.25) {\hw{T3: commit khi T2 valid}\\acc 26, iteration 4, cut, done};
\draw[wire] (sum)--(t3);
\draw[wire] (t3.east)--(7.45,-14.25) |- (m.east) node[pos=.61,tag,left]{acc:26};
\node[pin] (out) at (-2.45,-15.45) {\sig{acc[25:0]}, \sig{iterations\_done}, \sig{approx\_cut}, \sig{core\_done}};
\draw[wire] (t3)--(out);
}



\subsection{SAC, SBM và tích lũy không mất bit}

SAC dùng mask còn lại để tìm bit 1 có trọng số cao nhất, tạo
offset và mask cho bước tiếp theo. T1 chốt offset tích lũy,
token first/last, số vòng và cờ cắt. Token last được tạo khi
mask cạn hoặc khi đã dùng hết số vòng. Nếu không xử lý bit
fraction nào, token init-only chuyển tới T2 với term bằng 0.
Nhờ vậy đường dữ liệu commit vẫn chung với trường hợp có vòng.

T2 dịch phải cơ số \texttt{y\_base} theo offset của token.
Các vị trí bit thấp được đệm 0 từ trước, nên phép dịch từ
1 đến 12 giữ đủ số hạng trong lưới Q24. T3 chọn cơ số Y ở token
first hoặc accumulator ở các token sau, cộng term và lưu kết
quả. \texttt{core\_done} cùng có hiệu lực với kết quả của token
last đã commit. Số vòng fraction được cập nhật từ metadata của
token T2, không đếm slot init-only.

\begin{engineeringtable}{Trace số học bộ dữ liệu mẫu FT-05 trên MAC-F ở cấu hình OPS1/n12, $C=0$.}{tab:ft05-terms}
\begin{tabularx}{\linewidth}{@{}r l r Y@{}}
\toprule
Cạnh & Mốc & $M$ quan sát & Diễn giải \\
\midrule
0 & Accept & -- & Chốt A/B và cấu hình. \\
2 & \texttt{pair\_retire} & -- & Context core và C được chốt. \\
5 & T2 valid, vòng 1 & -- & T3 sắp commit số hạng đầu. \\
6 & T2 valid, vòng 2 & 23.633.920 & Cơ số 18.907.136 cộng term dịch phải 2. \\
7 & \texttt{core\_done} & 24.224.768 & Cộng thêm 590.848 từ dịch phải 5, $b=-11$, $t=2$. \\
14 & \texttt{out\_valid} & 24.224.768 & $D=\texttt{1ae34800}$, flags \texttt{00}. \\
\bottomrule
\end{tabularx}
\end{engineeringtable}

Ở n=0 với cùng A/B, init-only cho $M=18\,907\,136$, $t=0$ và
cờ còn số hạng được đặt. Mã đo được là \texttt{1a410000} với
flags \texttt{03}. Ở n12, X cạn sau hai vòng. Hai trường hợp
cho thấy thời gian xử lý phụ thuộc số bit 1 được chọn và không
đồng nhất với giá trị giới hạn cấu hình. Do bộ nhân luôn chạy ít nhất một chu kỳ ở pha lặp \texttt{T2/T3} sau \texttt{T1}, số chu kỳ thực thi của khối nhân lặp bằng $\max(1,t)$ chu kỳ cộng pha khởi tạo cố định, quyết định trực tiếp thành phần biến thiên trong độ trễ từng phép MAC.

\subsection{Adapter ở biên trước bộ cộng}

Nếu $M[25]=1$, significand tích nằm từ 2 đến dưới 4, nên scale
chuẩn hóa là $b+1$ và fraction lấy từ \texttt{\{M[24:0],1'b0\}}.
Nếu bit 25 bằng 0, scale giữ $b$ và fraction là
\texttt{\{M[23:0],2'b0\}}. Wrapper tính trước $b+1$ khi chốt
context, để adapter chọn giữa hai scale thay vì thêm carry-chain
vào đường từ core đến S1. Cả hai cách ánh xạ giữ đủ bit có giá
trị trong tích Q24. NF sử dụng biểu diễn đã chuẩn hóa để qua
Packer tích và Parser tích trước khi cộng, còn F dùng trực tiếp.
